\documentclass[14pt,a4paper]{extarticle}

\usepackage{geometry}
\geometry{left=30mm,right=15mm,top=20mm,bottom=20mm}
\usepackage{fontspec}
\usepackage{polyglossia}
\setmainlanguage{russian}
\setotherlanguage{english}
\setmainfont{Times New Roman}
\newfontfamily\cyrillicfont{Times New Roman}
\usepackage{indentfirst}
\usepackage{amsmath,amssymb}
\usepackage{csquotes}
\usepackage[
    backend=biber,
    bibencoding=utf8,
    sorting=none,
    style=gost-numeric,
    language=autobib,
    autolang=other,
    clearlang=true,
    defernumbers=false,
    sortcites=true,
    doi=false,
    isbn=false,
]{biblatex}
\addbibresource{/Users/gk/Code/super-duper-disser/itmo-phd-thesis-template-en/biblio/references.bib}
\addbibresource{/Users/gk/Code/super-duper-disser/itmo-phd-thesis-template-en/biblio/own.bib}

\setlength{\parindent}{1.25cm}
\setlength{\parskip}{0pt}
\linespread{1.15}
\emergencystretch=3em
\sloppy

\begin{document}

\begin{center}
\textbf{Введение}
\end{center}

\paragraph*{Актуальность темы исследования.}

Российская система территориального планирования связывает доступность объектов обслуживания с показателями национальных программ и проектов. Национальная цель по достижению комфортной и безопасной среды для жизни, поставлена в указе Президента РФ №309 от 7 мая 2024 года \cite{rfDecree309_2024}. Национальные направления развития выражаются в федеральных проектах, в частности, в проекте по формированию комфортной городской среды \cite{federalProjectFkgsPassport}. Прикладную часть программы федерального уровня реализовывают региональная программы, которые задают уже тот уровень, на котором локальный исполнительный орган власти непосредственно отвечает за достижение целевых показателей  эффективности \cite{spbRegionalFkgs}.

Несоответствие заключается в том, что оценки доступности и обеспеченности по нормативам отвечают на вопрос, где и что должно быть размещено, а ПКРТИ отвечает на вопрос, как развивать дорожно-транспортную инфраструктуру. В случае, если эти решения готовятся последовательно, возникает риск создания неэффективного плана застройки территории объектами обслуживания населения, поскольку согласование транспортного плана будет происходить на последующих этапах.

Задача обеспечения доступности объектов предложения не нова. Изначально, города строились с ориентацией на пешеходной движение, затем, с появлением всё новых видов мобильности, городская топология подстраивалась под них. Так, с увеличением расстояний и спроса на личный транспорт, городское пространство стало менее человеко-ориентированным, и, как следствие, менее доступным.

Таким образом, совместное развитие транспорта и застройки становится всё более актуальным. Подход транзитно-ориентированного развития (TOD) делает эту координацию более предметной в пространстве, определяя фокус на развитие территорий вокруг станций и коридоров общественного транспорта \cite{worldbankTod3v}. Историческая реконструкция TOD показывает, что этот подход возник как попытка связать землепользование и общественный транспорт в единую модель развития \cite{carlton2009tod}. Однако стоит отметить, что в модели TOD развитие сконцентрировано вокруг транспортных кластеров -- больших станций метро, вокзалов. В таком случае, зона влияния ограничена пешеходной изохроной вокруг рассматриваемых транспортных узлов, что не обеспечивает достаточной гибкости для развития на уровне всего города, ведь часто эти территории уже застроены. Доступность определяется как совместная характеристика застройки и транспорта.

Транспортная доступность не является постоянной характеристикой территории. В климатически зависимых регионах доступ к услугам меняется во времени: транспортные связи могут становиться недоступными, временно восстанавливаться или менять свою роль в сети. В арктических системах это особенно заметно, поскольку одни климатические процессы ухудшают работу наземной инфраструктуры, а другие временно создают новые возможности, например зимники и ледовые переправы. Поэтому задача доступности становится задачей оценки состояния сети в момент времени, а не только задачей расчета расстояния между поселением и ближайшим объектом обслуживания.

\paragraph*{Степень разработанности проблемы.}

Классическая линия моделей оптимального размещения объектов формализовала обслуживание как выбор узлов или площадок при заданном спросе и заданной сети расстояний. Задача размещения складов задала базовую постановку выбора складов с учетом затрат обслуживания спроса \cite{kuehn1963warehouse}. P-медиана и p-центр закрепили две ключевые логики: минимизацию суммарного расстояния и ограничение худшего расстояния до объекта \cite{hakimi1964medians}. Для общественных сервисов важным развитием стали покрывающие постановки. Покрытие множества минимизирует число объектов, достаточных для покрытия спроса в пределах нормативного расстояния \cite{toregas1971emergency}. Максимальное покрытие максимизирует покрытый спрос при ограниченном числе объектов \cite{church1974mclp}. В этих и других моделях оптимального размещения транспортная сеть обычно выражена фиксированной матрицей расстояний или времен, а не как самостоятельный объект решения.

Отдельная линия развития обеспечения доступности -- проектирование транспортной сети и проектирование сети с фиксированными затратами. Проектирование сети само по себе является классом оптимизационных задач \cite{magnanti1984networkDesign}. Проектирование сети с фиксированными затратами добавляет к выбору связей дискретные затраты открытия или изменения в сети \cite{powell1989networkDesign}. Алгоритмические работы по двойственному подъему показывают, что такая постановка требует специальных методов решения, а не сводится к обычному кратчайшему пути \cite{balakrishnan1989dualAscent}.

На стыке размещения и движения появились постановки, где результат зависит от качественных и количественных характеристик передвижения -- маршрутизации, времени прибытия и обслуживания потоков. В задачах размещения экстренных служб время движения и доступность объектов рассматриваются вероятностно \cite{daskin1987emergency}. Задача размещения коммивояжера связывает выбор местоположения с длиной маршрута обслуживания \cite{simchi1988tspLocation}. Задача размещения и маршрутизации объединяет размещение объектов и построение маршрутов в одной постановке \cite{laporte1988locationRouting}. Комбинированная постановка размещения и маршрутизации является отдельным направлением между размещением объектов и маршрутизацией транспорта \cite{min1998locationRouting}. Из этих работ видно, что маршрутная структура сети может менять доступность и перераспределять обслуживаемый спрос, как следствие, влияя на среднюю доступность в сети.

Дальнейшее развитие связано с переходом от сети как постоянной к сети как к части проектного решения. Модификация сети в задаче размещения показывает, что изменение связей может влиять на оптимальное размещение объектов \cite{berman1992networkModification}. Расширенное рассмотрение этой задачи связывает выбор объектов с топологией и структурой сети \cite{hodgson1992networkLocation}. Также было показано, что пространственная структура спроса в сети влияет на устойчивость результата размещения \cite{peeters1995spatialStructure}. Выделение пересадочных узлов и хабов также является самостоятельной группой методов оптимального размещения \cite{campbell1994hub}, находящих свое применение в задачах логистики \cite{klincewicz1998hub} \cite{revelle1996plantLocation}.

Класс задач совместного размещения объектов и проектирования сети включает в себя одновременный выбор мест размещения объектов обслуживания и изменения транспортных связей  \cite{daskin1993towardIntegrated} \cite{melkote1996integrated}. Решение задач покрытия в такой постановке показывает, что удовлетворение спроса зависит от изменений в сети \cite{melkote1998maximumCovering}. Решение этой задачи было также рассмотрено с учетом ограничения на вместимость входящих потоков и различной стоимостью добавления ребер в сети \cite{melkote2000capacitated}. Таким образом, совместное рассмотрение задачи оптимального размещения объектов и моделирования изменений сети был выделено в отдельную категорию Facility-location Network Design problem (FLND)\cite{melkote2001integrated}. Необходимо отметить, что при применении оптимизации размещения объектов к задаче покрытия спроса на доступ к социальной инфраструктуре в большинстве случаев необходимо рассматривать расширенную постановку задачи Facility location с учетом вместимости объектов обслуживания и размера спроса в заданной точке пространства.

Задача проектирования маршрутной сети (TNDP) формулируется как выбор набора линий или маршрутов на городском графе при заданном спросе и ограничениях связности \cite{mandl1980tndp}. Современные постановки планирования линий рассматривают последовательные этапы генерации линий, выбора линий и настройки частот \cite{schmidt2024linePlanning}. Обзоры TNDP показывают, что задача остается комбинаторно сложной и обычно балансирует пассажирские и операторские издержки \cite{duranMicco2022tndpReview}. Практические модели планирования общественного транспорта также используют ограничения длины маршрутов, пересадок и покрытия спроса \cite{sahin2023linePlanning}.

Меры снижения пространственного неравенства можно разделить на локальные, связующие и административные. Локальные меры добавляют новые сервисы в недостаточно обеспеченные территории. Связующие меры улучшают транспортную доступность между территориями. Административные меры меняют состав и логику управления агломерацией. Для задач обеспечения транспортной доступности особенно важна пара ``локальные -- связующие меры'', потому что она соответствует выбору между размещением объекта и изменением транспортной сети. Следовательно, задача обеспечения доступности не сводится к поиску новых площадок, а требует сравнения объектовых и сетевых вмешательств в единой модели.

\paragraph*{Цель исследования.}
Увеличение эффективности размещения объектов обслуживания населения в рамках процесса территориального планирования

\paragraph*{Задачи исследования.}
\begin{enumerate}
    \item Исследование системы управления территориальным планированием в задаче размещения объектов социального обслуживания населения;
    \item Разработка метода оценки транспортной доступности на основе критерия сезонной устойчивости;
    \item Разработка метода оптимизации сети в задаче оптимального размещения;
    \item Применение методов на примере агломераций Арктической зоны РФ.
\end{enumerate}

\paragraph*{Объект исследования.}
Подсистема территориального планирования в составе системы управления развитием территорий

\paragraph*{Предмет исследования.}
Методы и алгоритмы поддержки принятия решений по обеспечению транспортной доступности объектов обслуживания

\paragraph*{Методы исследования.}

Метод оценки транспортной доступности на основе критерия сезонной устойчивости, отличающийся применением методов сетевого анализа, обеспечивающий учет стохастической динамики многоуровневой сети.

Метод интеллектуальной поддержки принятия решений оптимального размещения объектов обслуживания, отличающийся применением оптимизации сети, обеспечивающий уменьшение необходимого количества объектов обслуживания.

\paragraph*{Научная новизна.}

В качестве частного случая этой общей схемы в данной главе используется сеть, заданная внешней средой (EFN), из арктической статьи. В этой модели явно выделяются два параметра: внешний параметр среды \(T\) и внутренний параметр связности графа \(G(t)\).

Объединяются изменения транспортной связности и размещение дополнительных сервисов, которые ранее рассматривались отдельно. Проверяемая гипотеза: улучшение связности между районами не более чем на 40\% может уменьшить минимальное число новых поликлиник, которые необходимо оптимально разместить \cite{kontsevik2024connectivityPlacement}.

Маршрутное проектирование рассматривается как Transit Network Design Problem (TNDP): требуется определить допустимый набор линий общественного транспорта, обеспечивающий удобные поездки для пассажиров при минимизации операторских затрат \cite{mandl1980tndp,duranMicco2022tndpReview,sahin2023linePlanning,morozov2026oneRuleTndp}.

\paragraph*{Теоретическая значимость.}

Теоретическая значимость состоит в развитии теории оценки транспортной доступности за счет учета влияния внешних факторов, качества составляющих сети, учета её конфигурации и учета эффекта взаимосвязи между слоями сети. За счет предложенных расширений получается учесть как изменение состояний самой сети, так и оценить эффект от направленного изменения сети.

\paragraph*{Практическая значимость.}
Практическая значимость работы состоит в программной реализации всех рассматриваемых методов, апробации методов на множестве территорий, что доказывает воспроизводимость вычислений и возможность применения разработанных методов для различных практических приложений.

\paragraph*{Положения, выносимые на защиту.}
\begin{enumerate}
    \item Метод оценки транспортной доступности на основе критерия сезонной устойчивости, отличающийся применением методов сетевого анализа, обеспечивающий учет стохастической динамики многоуровневой сети. Положение соответствует пункту 2 паспорта специальности.
    \item Метод интеллектуальной поддержки принятия решений оптимального размещения объектов обслуживания, отличающийся применением оптимизации сети, обеспечивающий уменьшение необходимого количества объектов обслуживания. Положение соответствует пункту 9 паспорта специальности.
\end{enumerate}

\paragraph*{Достоверность результатов.}
Степень достоверности научных достижений подтверждается корректным использованием методов, обоснованием постановки задач, экспериментальными исследованиями на реальных массивах данных, покрывающими разработанные технологии и алгоритмы.

\paragraph*{Апробация результатов.}
Результаты исследований обсуждались на конференциях:
\begin{enumerate}
    \item 40-я ежегодная конференция AAAI по искусственному интеллекту (2026) --- 2-й семинар по ИИ для городского планирования, 2026, рейтинг CORE: A*, стендовый и устный доклад;
    \item Научная конференция AMS (2026) в области городского планирования, специальная секция, устный доклад;
    \item Международная школа и конференция по науке о сетях (NetSci 2026), основная программа, стендовый доклад;
    \item Международная конференция по вычислительной науке и ее приложениям (ICCSA 2023/2024), устный доклад.
\end{enumerate}

\paragraph*{Личный вклад автора.}
В работах, выполненных в соавторстве, вклад автора заключается в построении моделей, разработке и реализации методов и алгоритмов, разработке, проведении вычислительных экспериментов и подготовке обзоров литературы.

\paragraph*{Структура и объем диссертации.}
Во введении сформулированы цель и задачи исследования, обоснованы актуальность и научная новизна работы. Также, во введении перечислены основные положения, выносимые на защиту диссертационной работы; представлены практическая и теоретическая значимость.

В первой главе предоставлен литературный обзор по теме диссертации. Показаны состояние предметной области, а также недостатки существующих подходов.

Во второй главе описан метод оценки транспортной доступности на основе критерия сезонной устойчивости, описаны результаты применения методов сетевого анализа, показаны результаты учета динамики сезонного изменения сети.

В третьей главе описан метод интеллектуальной поддержки принятия решений размещения объектов обслуживания; описаны основные отличия оптимизации размещения с учетом изменчивости сети.

В четвертой главе описаны экспериментальные исследования эффективности предложенных методов для прикладных задач управления в системе управления развитием территории.

\paragraph*{Соответствие специальности 2.3.4 (в порядке степени соответствия).}

\begin{itemize}
    \item[2.] Разработка математических моделей и критериев эффективности, качества и надёжности организационных систем;
    \item[9.] Разработка методов и алгоритмов интеллектуальной поддержки принятия управленческих решений в организационных системах.
\end{itemize}

\end{document}
